% Part: incompleteness
% Chapter: theories-computability
% Section: intepretability

\documentclass[../../../include/open-logic-section]{subfiles}

\begin{document}

\olfileid{inc}{tcp}{itp}

\olsection{$\Th{Q}$નું અર્થઘટન થઈ શકે એવા સિદ્ધાંતો અનિર્ણેય છે}

આ પરિણામોને વધુ મજબૂત
બનાવી શકીએ. અનૌપચારિક રીતે,
એક ભાષા~$\Lang{L_1}$નું
બીજી ભાષા~$\Lang{L_2}$માં
અર્થઘટન એટલે
$\Lang{L_1}$ના વિશ્વ,
સંબંધ-ચિહ્નો અને
વિધેય-ચિહ્નોને
$\Lang{L_2}$નાં
!!{formula} વડે વ્યાખ્યાયિત
કરવાં. અહીં એટલો સમય
લેતા નથી, છતાં આ
વ્યાખ્યાને ચોક્કસ બનાવી
શકાય છે.

\begin{thm}
  માનો કે $\Th{T}$ એવી
  ભાષાનો સિદ્ધાંત છે જેમાં
  અંકગણિતની ભાષાનું એવું
  અર્થઘટન શક્ય છે કે
  $\Th{T}$ એ $\Th{Q}$ના
  અર્થઘટન સાથે સુસંગત છે.
  તો $\Th{T}$ અનિર્ણેય છે.
  જો $\Th{T}$ એ
  $\Th{Q}$ની સ્વયંસિદ્ધિઓના
  અર્થઘટનને સાબિત કરે,
  તો $\Th{T}$નું કોઈ
  સુસંગત વિસ્તરણ નિર્ણેય
  નથી.
\end{thm}

સાબિતી અગાઉના પ્રમેયની
સાબિતીમાં નાનો ફેરફાર
કરીને મળે છે: વિરોધી
ઉદાહરણ હોય તો તેનાથી
$\Th{Q}$ અને
$\Th{\bar Q}$ને અલગ
પાડી શકીએ. પસંદગીના
સ્વયંસિદ્ધિ સહિતનો
ઝર્મેલો--ફ્રેન્કલ ગણસિદ્ધાંત
$\Th{ZFC}$ સામાન્ય ગણિતનો
મોટો ભાગ વિકસાવી શકાય
એવો શક્તિશાળી
સ્વયંસિદ્ધીય આધાર છે.
$\Th{ZFC}$માં પ્રાકૃતિક
સંખ્યાઓ વ્યાખ્યાયિત કરી
શકાય છે, અને આ અર્થઘટન
હેઠળ $\Th{Q}$ની
સ્વયંસિદ્ધિઓ સાચી છે.
આથી આપણને મળે છે:

\begin{cor}
$\Th{ZFC}$નું કોઈ
સુસંગત વિસ્તરણ નિર્ણેય નથી.
\sourcecorrection{OLINC-042}{સ્થિર અંગ્રેજી
મૂળના અનુવિધાનમાં
``સુસંગત'' શરત છૂટી છે.
અસુસંગત વિસ્તરણ દરેક
વાક્ય સાબિત કરે છે અને
તેથી નિર્ણેય હોય છે.
ઉપરના પ્રમેય પરથી મળતા
અનુવિધાનમાં સુસંગતતા
જરૂરી છે.}
\end{cor}

\begin{cor}
$\Th{ZFC}$નું પૂર્ણ,
સુસંગત અને સંગણનીય
રીતે !!{axiomatizable}
વિસ્તરણ હોતું નથી.
\end{cor}

$\Th{ZFC}$ની ભાષામાં
માત્ર એક દ્વિચલ સંબંધ
$\in$ છે. (હકીકતમાં
સમતા પણ જરૂરી નથી.)
આથી મળે છે:

\begin{cor}
દ્વિચલ સંબંધ-ચિહ્ન
ધરાવતી કોઈપણ ભાષાનો
પ્રથમ-ક્રમનો તર્ક
અનિર્ણેય છે.
\end{cor}

આ પરિણામ બે એકચલ
વિધેય-ચિહ્ન ધરાવતી કોઈપણ
ભાષા માટે પણ લાગુ પડે છે,
કારણ કે તેમની મદદથી
દ્વિચલ સંબંધ-ચિહ્નનું
અનુકરણ કરી શકાય છે.
હમણાં આપેલાં પરિણામો
ચોક્કસ સીમાએ છે: માત્ર
\emph{એકચલ} સંબંધ-ચિહ્નો
અને વધુમાં વધુ એક
\emph{એકચલ} વિધેય-ચિહ્ન
ધરાવતી ભાષાનો પ્રથમ-ક્રમનો
તર્ક નિર્ણેય છે.

એક વધારાની માહિતી.
પ્રમાણભૂત નમૂનામાં સાચાં
હોય એવાં, ભાષા
$\Obj 0$, $'$, $+$,
$\times$, $<$નાં વાક્યોનો
ગણ અનિર્ણેય છે તે આપણે
જાણીએ છીએ. હકીકતમાં
$<$ને બાકીનાં ચિહ્નો
વડે વ્યાખ્યાયિત કરી
શકાય છે, અને પછી
$+$ને $\times$ અને $'$
વડે વ્યાખ્યાયિત કરી
શકાય છે. તેથી ભાષા
$\Obj 0$, $'$, $\times$નાં
સાચાં વાક્યોનો ગણ
અનિર્ણેય છે. બીજી બાજુ,
પ્રેસબર્ગરે બતાવ્યું કે
અંકગણિતની ભાષામાં
સાચાં હોય એવાં, ભાષા
$\Obj 0$, $'$, $+$નાં
વાક્યોનો ગણ નિર્ણેય છે.
જોકે આ નિર્ણય-પ્રક્રિયા
ગણતરી માટે વ્યવહારુ નથી.

\end{document}
